% !TEX root = 00_main_aaai_2027.tex
\section{Evaluate Conditional Probability Forecasts}
\label{sec:proposal}
For ratings $r\in\{0,\ldots,5\}$ and input $x=(y,y^*)$, we evaluate a judge $q_\theta(\cdot\mid x)$ with a bounded, strictly proper score matched to the label encoding~\citep{gneiting2007strictly}.
The \emph{normalized Brier score} on the one-hot encoding and the ranked probability score (RPS) on the cumulative encoding are
\begin{align}
L_{\rm Brier}(q,r)&=\frac12\sum_{j=0}^{5}(q_j-\mathbf{1}\{r=j\})^2,\label{eq:brier-norm}\\
L_{\rm RPS}(q,r)&=\frac{1}{5}\sum_{j=0}^{4}(Q_j-\mathbf{1}\{r\leq j\})^2,
\end{align}
where $Q_j=\sum_{k=0}^j q_k$.
On a hard forecast $\delta_a$, normalized Brier equals $0$--$1$ error $\mathbf{1}\{a\neq r\}$ and RPS equals normalized absolute error $\lvert a-r\rvert/5$ (\appref{app:rps}{the RPS appendix} records the identities).
Both scores are unbiased estimates of their expected value with a single human rating per item.
The appendix on the ranked probability score explains RPS in detail: intuition, propriety and why it suits ordinal human ratings.
Their expected excess losses over the oracle are $\frac12\sum_j(q_j-p_j)^2$ and $\frac{1}{5}\sum_{j=0}^4(Q_j-P_j)^2$.
These vanish only when the categorical distributions agree.

Cross-entropy (CE), $\mathbb{E}_{x,R}[-\log q_\theta(R\mid x)]$, is a strictly proper but unbounded secondary diagnostic, and a convenient training loss.
It differs from the expected Kullback--Leibler (KL) divergence
\begin{equation}
\min_\theta\ \mathbb{E}_{x}\big[\KL(p(\cdot\mid x)\,\|\,q_\theta(\cdot\mid x))\big].
\label{eq:kl}
\end{equation}
by an entropy term independent of $\theta$.
The target $p$ is a population distribution, not a sparse empirical histogram.
One observed rating per item supplies a valid conditional log-loss observation.
It does not require marginal pooling and does not identify an unconstrained distribution at that item without shared structure or further assumptions.
CE has excess loss $\KL(p\|q)$, assuming forecast support covers the population support.

\paragraph{Skill and agreement strata.}
The per-sample risk of a score $S$ over held-out items $x$ with $k_x$ ratings each, $\hat R_S=\frac1{|D|}\sum_{x\in D}\frac1{k_x}\sum_{j=1}^{k_x}S\big(q(\cdot\mid x),r_{x,j}\big)$, is unbiased for the expected score for any number of ratings per item, including one.
The skill score $1-\hat R_S(\text{judge})/\hat R_S(\text{base rate})$, with the marginal rating distribution as base rate, is $0$ for a judge as good as the base rate, negative for a worse one, and at most $1$.
Being a ratio of two risk estimates, it is consistent rather than unbiased.
The risk ranks judges, but the distance to $p$ (the expected excess loss) is the risk minus the uncertainty term $\mathbb{E}_x S(p_x,p_x)$, for normalized Brier $\frac12(1-\sum_j p_j^2)$ and for each RPS threshold $P_j(1-P_j)$; these depend on how often two ratings of one item agree, so estimating the distance needs at least two ratings per item (finite statements Lean-checked with numerics, Experiment~53).
Items with at least four ratings are stratified by human agreement measured on one random half of their ratings (so at least two ratings define agreement), and every score is computed on the other half, so a stratum is not chosen with the ratings that score it.

Empirical evaluation averages losses over observed ratings, including with one rating per input.
A histogram of two or three ratings remains a noisy estimate, and pooling marginal histograms discards the input--rating association.
Probability fitting also does not remove finite-sample selection error, model misspecification, or optimization failure.
We retain Pearson, Spearman and Kendall as diagnostics, with normalized scalar error and signed bias when score levels matter.
The Brier columns in Table~\ref{tab:population} report the unnormalized sum of squared errors $\sum_j(q_j-\mathbf{1}\{r=j\})^2$, twice the normalized score in Eq.~\ref{eq:brier-norm}.
